\section{Threats to Validity}
\label{sec:threats}

\textbf{Internal validity.} We randomized the run order, but differences of a few percentage points between configurations remain within run-to-run variation, which reached 6.3 points for A3. Overhead differences smaller than this should not be interpreted. Configurations C and E classified and responded after each workload ended, so their decision and freeze times do not include the time an online detector would need. The model was trained on 16 windows from one calibration run per workload. The in-kernel window state is not updated atomically, so concurrent events on several CPUs can make the in-kernel counts differ slightly from the user-space counts. The in-kernel decision latency rests on 20 values, one per flagged run.

\textbf{Construct validity.} Throughput is measured in operations per workload, and an operation means different things in different workloads, so overheads are comparable within a workload but not across workloads. The transfer latency includes the time a record waits for the consumer, not only the time to cross the kernel boundary. The in-kernel decision latency includes the cost of recording the triggering event and of reserving the alert record. The user-space inference latency was measured in Python. The freeze latency was measured on containers whose workload had finished, and each freeze was timed with a polling loop that reads \texttt{cgroup.events} with a 0.1~ms sleep, so latencies near that interval carry corresponding resolution limits.

\textbf{External validity.} All runs used one host, one kernel version and the WSL2 virtualization layer, and the per-event cost may differ on bare metal or other hypervisors. The workloads are synthetic and short, run one at a time, and use Docker without an orchestrator. The probe hooks five system calls on entry. It does not see \texttt{execveat}, \texttt{open}, renames or the results of the calls it hooks, so it would miss behaviour that uses only those paths. The detector is a single linear model trained on two attack types, and we make no claim about real attacks, adaptive attackers or other model families. We did not compare against existing tools such as Falco or Tetragon.

\textbf{Conclusion validity.} Each cell has ten runs. We report means, standard deviations and ranges and avoid claims that depend on small differences. We did not apply significance tests, because the main effects, such as overheads above 40\% and the complete miss of A3, are far larger than the observed run-to-run variation, and the absence of a measurable cost for in-kernel inference is a statement about differences that are indistinguishable from that variation.

\section{Conclusion}
\label{sec:conclusion}

We evaluated an eBPF and machine-learning runtime-security pipeline for containers under one measurement protocol, holding telemetry, features, model and workloads constant across five configurations and 350 runs in randomized order. Five findings stand out. First, after we corrected a consumer defect, the telemetry was complete: 8.87 million events with no loss. Second, its cost was a per-event cost of 12 to 25~\textmu s on our testbed, so overhead ranged from nothing measurable to 57\%, depending on how many monitored system calls a workload made. Third, running the same model inside the probe in Q16.16 arithmetic added no measurable overhead, changed no decision and made the decision available in user space about 20~\textmu s after the triggering syscall, instead of after a full window. Fourth, a detector that was perfect on its trained anomalies missed every run of an unseen connection burst, because it had learned that connections indicate benign activity. Fifth, containment through the cgroup v2 freezer took effect in well under a millisecond and was always reversible.

These results suggest that, for eBPF-based ML, the value of kernel placement lies in incremental, per-event decisions inside the event path, not in faster arithmetic. Future work should act on in-kernel alerts directly so that containment reaches running workloads within microseconds, enlarge the training set, add realistic benign workloads and real attack tools, compare against existing runtime-security tools, and repeat the study on bare-metal hosts and under Kubernetes.

\section*{Data Availability}

\begingroup\raggedright
The code, raw run records, processed tables and figures are available at \url{https://github.com/SujitSaiY2007/Computer-Networks-Case-Study-}. The results in this paper come from commit \texttt{d0870ab}. The run records are in \path{experiments/raw}, and the run order is in \path{experiments/manifests/run_order.json}. The datasets produced before each correction are kept, with their calibration data, in \path{experiments/raw_v1_early_exit_loader}, \path{experiments/raw_v2_same_code_anomalies} and \path{experiments/raw_v3_python_configD}, so that the effect of each change can be inspected.
\par\endgroup
